\documentclass[11pt]{amsart}

\usepackage[T1]{fontenc}
\usepackage{lmodern}
\usepackage{microtype}
\usepackage{amsmath,amssymb,amsthm}
\usepackage[hidelinks]{hyperref}

\newtheorem{theorem}{Theorem}

\newcommand{\I}{\mathcal I}

\title{A 79-Vertex Counterexample to Conjecture~4.1 of Hurlbert and Kamat}
\subjclass[2020]{05C69, 05D05}
\keywords{independent set, maximum star center, HK tree, tree,
Erd\H{o}s--Ko--Rado}
\date{}

\hypersetup{
  pdftitle={A 79-Vertex Counterexample to Conjecture 4.1 of Hurlbert and Kamat}
}

\begin{document}

\begin{abstract}
We exhibit a tree on $79$ vertices with no vertex of degree two such that,
for $r=18$, a degree-three vertex is contained in more independent $r$-sets
than any leaf. Thus the tree is not $18$-HK, disproving Conjecture~4.1 of
Hurlbert and Kamat. The comparison is certified by exact
independence-polynomial identities.
\end{abstract}

\maketitle

\section{Introduction}

For a graph $G$ and $1\le r\le \alpha(G)$, let $\I^r(G)$ be the family of
independent $r$-subsets of $V(G)$, and let
\[
\I_v^r(G)=\{A\in\I^r(G):v\in A\}.
\]
A vertex $v$ is a \emph{maximum $r$-center} if $|\I_v^r(G)|$ is maximum over
$V(G)$. Following Estrugo and Pastine~\cite{EstrugoPastine}, a tree is
$r$-HK if some leaf is a maximum $r$-center, and is HK if it is $r$-HK for
every admissible $r$.

Counterexamples to the unrestricted assertion that every tree is HK are
known~\cite{Borg,FeghaliJohnsonThomas}. Hurlbert and Kamat subsequently
conjectured that every tree with no vertex of degree two is
HK~\cite[Conjecture~4.1]{HurlbertKamat}; see also
\cite[Section~4]{HurlbertSurvey}. We give a counterexample.

\begin{theorem}\label{thm:main}
There is a tree $T$ on $79$ vertices, with no vertex of degree two, and a
vertex $c$ of degree three such that
\[
|\I_c^{18}(T)|
=1{,}526{,}033{,}336{,}050{,}350
>
\max_{u\text{ a leaf of }T}|\I_u^{18}(T)|
=1{,}524{,}522{,}945{,}307{,}798.
\]
Consequently, $T$ is not $18$-HK, and Conjecture~4.1 of Hurlbert and Kamat
is false.
\end{theorem}

\section{Construction and verification}

\begin{proof}[Proof of Theorem~\ref{thm:main}]
Let $t_1=9$ and $t_2=t_3=8$. Define $T$ by
\[
\begin{aligned}
V(T)={}&\{c\}\cup\{h_i:1\le i\le3\}
\cup\{s_{ij}:1\le i\le3,\ 1\le j\le t_i\}\\
&\cup\{\ell_{ijk}:1\le i\le3,\ 1\le j\le t_i,\ k\in\{1,2\}\},
\end{aligned}
\]
and with edges
\[
ch_i,\qquad h_is_{ij},\qquad s_{ij}\ell_{ijk}
\]
for all admissible indices. Every vertex other than $c$ has a unique
neighbor closer to $c$, so $T$ is a tree. Moreover,
\[
|V(T)|=1+3+3(t_1+t_2+t_3)=79.
\]
Also,
\[
\deg(c)=\deg(s_{ij})=3,\qquad
\deg(h_i)=t_i+1\in\{9,10\},\qquad
\deg(\ell_{ijk})=1,
\]
so $T$ has no vertex of degree two. The three vertices $h_i$ together with
all $50$ leaves form an independent set, so $18\le\alpha(T)$.

For a graph $F$, write
\[
P_F(z)=\sum_{k\ge0}i_k(F)z^k,
\]
where $i_k(F)$ is the number of independent $k$-sets of $F$. For every graph
$G$ and vertex $v\in V(G)$,
\begin{equation}\label{eq:star}
|\I_v^r(G)|=[z^{r-1}]P_{G-N_G[v]}(z),
\end{equation}
where $N_G[v]$ is the closed neighborhood of $v$.

A support vertex $s_{ij}$ together with its two leaf neighbors induces a
three-vertex path, whose independence polynomial is
\[
S(z)=1+3z+z^2.
\]
For $t\ge0$, set
\[
H_t(z)=S(z)^t+z(1+z)^{2t}.
\]
This is the independence polynomial of a branch with one hub, $t$ support
vertices, and two leaves adjacent to each support vertex: the first term
omits the hub, and the second includes it.

Deleting $N_T[c]$ leaves $25$ disjoint three-vertex paths, and hence
\begin{equation}\label{eq:center}
P_{T-N_T[c]}(z)=S(z)^{25}.
\end{equation}
The leaves form two orbits. Let $u_8$ be a leaf in a branch with eight
support vertices, and let $u_9$ be a leaf in the branch with nine support
vertices. After deleting a leaf and its support vertex, its sibling remains
isolated. Splitting the remaining graph according to whether $c$ is chosen
gives
\begin{align}
P_{T-N_T[u_8]}(z)
  &=(1+z)\bigl(zS(z)^{24}+H_9(z)H_8(z)H_7(z)\bigr),
    \label{eq:leaf8}\\
P_{T-N_T[u_9]}(z)
  &=(1+z)\bigl(zS(z)^{24}+H_8(z)^3\bigr).
    \label{eq:leaf9}
\end{align}

For nonnegative integers $q,m$, with the convention $\binom{n}{k}=0$ for
$k<0$ or $k>n$, exact coefficients are obtained from
\begin{equation}\label{eq:coeff}
[z^m]S(z)^q
=
\sum_{j=\max\{0,m-q\}}^{\lfloor m/2\rfloor}
\binom{q}{j}\binom{q-j}{m-2j}3^{m-2j}
\end{equation}
and
\[
[z^m]H_t(z)=[z^m]S(z)^t+\binom{2t}{m-1}.
\]
Applying \eqref{eq:coeff} and ordinary coefficient convolution to
\eqref{eq:center}--\eqref{eq:leaf9} yields
\begin{align*}
[z^{17}]P_{T-N_T[c]}(z)
  &=1{,}526{,}033{,}336{,}050{,}350,\\
[z^{17}]P_{T-N_T[u_8]}(z)
  &=1{,}524{,}522{,}945{,}307{,}798,\\
[z^{17}]P_{T-N_T[u_9]}(z)
  &=1{,}521{,}164{,}183{,}552{,}880.
\end{align*}
In particular,
\[
|\I_c^{18}(T)|-|\I_{u_8}^{18}(T)|
=1{,}510{,}390{,}742{,}552>0.
\]
By \eqref{eq:star}, the star centered at $c$ is larger than every
leaf-centered $18$-star. Hence no leaf is a maximum $18$-center, and $T$ is
not $18$-HK.
\end{proof}

\medskip
\noindent\emph{Remark.} No minimality claim is made for the order of the
construction.

\begin{thebibliography}{9}

\bibitem{Borg}
P.~Borg,
\emph{Stars on trees},
Discrete Math. \textbf{340} (2017), no.~5, 1046--1049.
\href{https://doi.org/10.1016/j.disc.2016.11.002}
{doi:10.1016/j.disc.2016.11.002}.

\bibitem{EstrugoPastine}
E.~J.~J. Estrugo and A.~Pastine,
\emph{On stars in caterpillars and lobsters},
Discrete Appl. Math. \textbf{298} (2021), 50--55.
\href{https://doi.org/10.1016/j.dam.2021.03.013}
{doi:10.1016/j.dam.2021.03.013}.

\bibitem{FeghaliJohnsonThomas}
C.~Feghali, M.~Johnson, and D.~Thomas,
\emph{Erd\H{o}s--Ko--Rado theorems for a family of trees},
Discrete Appl. Math. \textbf{236} (2018), 464--471.
\href{https://doi.org/10.1016/j.dam.2017.10.009}
{doi:10.1016/j.dam.2017.10.009}.

\bibitem{HurlbertSurvey}
G.~Hurlbert,
\emph{A survey of the Holroyd--Talbot conjecture},
in G.~O.~H. Katona, B.~Patk\'os, and C.~Tompkins (eds.),
\emph{Sum(m)it280: Surveys in Extremal Combinatorics and Combinatorial
Geometry}, Bolyai Society Mathematical Studies, vol.~32,
Springer, Cham, 2026, pp.~183--201.
\href{https://doi.org/10.1007/978-3-032-18810-6_10}
{doi:10.1007/978-3-032-18810-6\_10}.

\bibitem{HurlbertKamat}
G.~Hurlbert and V.~Kamat,
\emph{On intersecting families of independent sets in trees},
Discrete Appl. Math. \textbf{321} (2022), 4--9.
\href{https://doi.org/10.1016/j.dam.2022.06.028}
{doi:10.1016/j.dam.2022.06.028}.

\end{thebibliography}

\end{document}